\documentclass[level=1]{kidscode}
%% \gl{term}{where you meet it}{meaning}{example (optional)}
\newcommand{\gl}[4]{\par\noindent\needspace{3\baselineskip}{\bfseries #1}\hfill{\footnotesize\color{kcgray}#2}\par\nopagebreak
  \noindent #3\ifstrempty{#4}{}{\par\nopagebreak\noindent{\small\color{kcgray!80!black}Example: #4}}\par\vspace{5pt}}
\newcommand{\letter}[1]{\par\needspace{7\baselineskip}\vspace{4pt}\noindent{\Large\bfseries\color{kcaccent}#1}\par\nopagebreak\vspace{2pt}}
\begin{document}
\kctitle{Glossary}{Every new word in Level 1, in A to Z order}

\begin{tip}[How to use this glossary]
Words are in alphabetical order. The gray label tells you which week you first meet the word
(W1 = week 1, B1 = buffer session 1). Words in \texttt{this font} are typed exactly like that in Python.
Found a word that is not here? Write it on the last page with your own meaning!
\end{tip}

\begin{multicols}{2}
\letter{A}
\gl{accessibility}{W5}{Making things that everyone can use, including people who see colors differently or need bigger text.}{Using color-blind-safe colors and also different shapes.}
\gl{algorithm}{W1}{A list of exact steps, in the right order, that solves a problem or does a job.}{A recipe for a sandwich.}
\gl{angle}{W2}{How far you turn, measured in degrees. A full turn is 360 degrees. A shape with $n$ sides turns $360 \div n$ at each corner.}{A square turns 90 degrees at each corner.}

\letter{B}
\gl{backdrop}{W1}{The picture behind everything on the Scratch stage.}{}
\gl{broadcast}{W5}{A Scratch message sent to every sprite. Sprites with a matching \enquote{when I receive} block start running.}{\enquote{broadcast start game}.}
\gl{bug}{W1}{A mistake in a program that makes it do something we did not want. Bugs are normal: every programmer meets them every day.}{The turtle turns 90 instead of 120.}

\letter{C}
\gl{call}{W10}{To use a function, by writing its name with parentheses.}{\texttt{house(0, 0, ORANGE)} calls the function \texttt{house}.}
\gl{color code}{W9}{A color written as \# and six letters or numbers, so the computer knows exactly which color you mean.}{\texttt{"\#0072B2"} is blue.}
\gl{comment}{W7}{A note for humans in the code. Python ignores it. In Python it starts with \texttt{\#}.}{\texttt{\# draw the roof}}
\gl{comparison}{W9}{A question that is True or False, made with $>$ (bigger), $<$ (smaller), \texttt{==} (equal) or \texttt{!=} (not equal).}{\texttt{stars > 5}}
\gl{conditional}{W3}{Code that only runs if something is true. Also called an \textbf{if}.}{\enquote{If touching the edge, then bounce.}}
\gl{coordinates}{W1}{Two numbers, $x$ and $y$, that give a place on the screen. $x$ is left--right, $y$ is down--up. The middle is $(0, 0)$.}{\texttt{t.goto(200, 150)}}
\gl{costume}{W2}{One of the pictures a Scratch sprite can wear. Switching costumes quickly makes animation.}{}
\gl{credit}{W4}{Saying who made the art, sounds or ideas you used. It is fair and kind.}{\enquote{Cat sounds from the Scratch library.}}

\letter{D}
\gl{debug}{W1}{To find a bug and fix it. A good strategy: read the error, check the last thing you changed, test one small part at a time.}{}
\gl{decomposition}{W5}{Breaking a big problem into smaller parts that are easier to solve.}{A house = walls + roof + door + windows.}
\gl{\texttt{def}}{W10}{The Python word that starts a new function (your own command).}{\texttt{def square(size):}}
\gl{driver}{W1}{In pair programming, the person using the mouse and keyboard.}{}

\letter{E}
\gl{\texttt{elif}}{W9}{Short for \enquote{else if}: another question to check when the first \texttt{if} was false.}{\texttt{elif n == 10:}}
\gl{else}{W3}{The part of an if that runs when the question is false.}{}
\gl{error message}{W7}{The red text Python shows when it cannot run your code. It tells you the line number and a clue.}{\texttt{NameError: name 'Print' is not defined}}
\gl{event}{W3}{Something that happens and starts code running.}{Clicking the green flag, pressing a key.}
\gl{extension}{W2}{Extra blocks you can add to Scratch, like the Pen.}{}

\letter{F}
\gl{feedback}{W3}{Helpful comments about someone's work. We use \enquote{two stars and a wish}: two things you like, one idea to improve.}{}
\gl{\texttt{for} loop}{W8}{Python's repeat block. The lines pushed in underneath it are repeated.}{\texttt{for i in range(4):}}
\gl{forever}{W2}{A Scratch loop that never stops (until you click the red stop sign).}{}
\gl{function}{W10}{A named set of steps you can use again and again. Python's version of a My Block.}{\texttt{square(50)}}

\letter{G}
\gl{green flag}{W1}{The Scratch button that starts a project.}{}

\letter{I}
\gl{if}{W3}{A block or line that asks a question and only runs its steps when the answer is true.}{\texttt{if mood == "happy":}}
\gl{indentation}{W8}{Spaces at the start of a line (4 in Python). They show which lines are inside a loop, an if or a function.}{}
\gl{input}{W9}{Information that goes into a program, often typed by the user. In Python \texttt{input()} asks a question and gives back the answer as text.}{\texttt{answer = input("How many? ")}}
\gl{integer (\texttt{int})}{W9}{A whole number, like 3, 0 or $-7$. \texttt{int()} turns text like \texttt{"5"} into the number 5.}{\texttt{n = int(answer)}}

\letter{L}
\gl{list}{W8}{Several values kept together in square brackets. Python counts the places from 0, so \texttt{PALETTE[0]} is the first one.}{\texttt{PALETTE = ["\#E69F00", "\#56B4E9", \ldots]}}
\gl{loop}{W2}{Code that repeats. Loops save you writing the same steps again and again.}{repeat 4, \texttt{for i in range(4):}}

\letter{M}
\gl{My Block}{W5}{A block you make yourself in Scratch with \enquote{Make a Block}. It can have inputs.}{\enquote{define flower (size)}}

\letter{N}
\gl{navigator}{W1}{In pair programming, the person who reads the instructions, thinks ahead and spots bugs.}{}
\gl{nested loop}{W8}{A loop inside another loop.}{Repeat 6 times: draw a square.}

\letter{O}
\gl{output}{W7}{What a program shows or does: text, a drawing, a sound.}{}

\letter{P}
\gl{pair programming}{W1}{Two people coding together on one computer: a driver and a navigator who swap roles.}{}
\gl{parameter}{W5}{An input that a My Block or function needs. It makes the block flexible.}{In \texttt{square(size)}, \texttt{size} is a parameter.}
\gl{pen}{W2}{Lets a sprite or turtle draw a line where it moves. Pen up: no line. Pen down: draw.}{}
\gl{pixel}{B1}{One tiny square of color. Screens make pictures out of millions of pixels.}{}
\gl{playtest}{W11}{Letting someone try your program so you can see what works and what confuses them.}{}
\gl{\texttt{print}}{W7}{The Python command that shows text in the Shell.}{\texttt{print("Hello!")}}
\gl{privacy}{W1}{Keeping personal information (real name, address, school, photos, passwords) safe. We use nicknames.}{}
\gl{program}{W1}{A set of instructions a computer follows.}{}
\gl{Python}{W6}{A programming language where you type words instead of dragging blocks.}{}

\letter{R}
\gl{random}{W4}{Picked by chance, so you cannot guess it. Great for surprises and games.}{\texttt{random.randint(1, 6)}, \enquote{pick random 1 to 6}}
\gl{\texttt{range}}{W8}{Tells a \texttt{for} loop how many times to repeat.}{\texttt{range(5)} repeats 5 times.}
\gl{remix}{W4}{Making something new from someone else's work, and giving them credit.}{}
\gl{repeat}{W2}{A Scratch loop that runs the blocks inside a set number of times.}{}
\gl{run}{W1}{To start a program. In Thonny, the green Run button or \keys{F5}.}{}

\letter{S}
\gl{script}{W1}{A stack of Scratch blocks that are joined together.}{}
\gl{sensing}{W3}{Scratch blocks that check things, like \enquote{touching edge?} or \enquote{key space pressed?}.}{}
\gl{sequence}{W1}{Steps that happen one after another, in order. Order matters!}{}
\gl{Shell}{W7}{The part of Thonny at the bottom where \texttt{print} shows text and you type answers.}{}
\gl{sprite}{W1}{A character or object in Scratch that you can program.}{The Scratch cat.}
\gl{stage}{W1}{The area in Scratch where the sprites move and the project plays.}{}
\gl{string}{W9}{Text in a program, written inside quotes. \texttt{input()} always gives a string.}{\texttt{"day"}, \texttt{"25"}}
\gl{syntax}{W7}{The spelling and grammar rules of a programming language.}{}
\gl{syntax error}{W7}{An error because the code breaks Python's rules, like a missing bracket or colon.}{\texttt{print("hi"}}

\letter{T}
\gl{test}{W7}{Running a program on purpose to check that it works. \texttt{check\_my\_work.py} runs tests for you.}{}
\gl{Thonny}{W7}{The app we use to write and run Python.}{}
\gl{turtle}{W7}{A Python drawing tool: an arrow you move around the screen that draws lines.}{\texttt{t = turtle.Turtle()}}

\letter{V}
\gl{variable}{W4}{A labeled box that remembers a value, like a score or an answer. The value can change.}{\texttt{score = 0}}
\end{multicols}

\subsection*{More words from your workbook}
\begin{multicols}{2}
\begin{description}[leftmargin=0pt,itemsep=3pt,font=\bfseries]
\item[accessible] Easy for everyone to use, including people who see or hear differently. {\small\color{kcgray}(week 5)}
\item[bug / debugging] A mistake in a program / finding and fixing it. {\small\color{kcgray}(week 1)}
\item[checkpoint] A short check that shows what you know and what to practice next. {\small\color{kcgray}(week 6)}
\item[debugging] Finding and fixing bugs, step by step, like a detective. {\small\color{kcgray}(week 6)}
\item[Demo] Short for demonstration: running your program and showing how it works. {\small\color{kcgray}(week 12)}
\item[driver / navigator] The driver clicks; the navigator reads, checks and helps. {\small\color{kcgray}(week 1)}
\item[Editor] The top part of Thonny where you write your program. {\small\color{kcgray}(week 7)}
\item[Feedback card] A card with two stars and a wish for an artist. {\small\color{kcgray}(week 12)}
\item[if / elif / else] Python's way to choose: \enquote{if this, else if that, otherwise\ldots} {\small\color{kcgray}(week 9)}
\item[if-else] Runs one set of blocks for yes and another for no. {\small\color{kcgray}(week 3)}
\item[Iterate] To improve something by going round the build--run--check loop again and again. {\small\color{kcgray}(week 11)}
\item[Milestone] A checkpoint on the way to finishing a project. {\small\color{kcgray}(week 11)}
\item[Off-by-one bug] When a loop runs one time too many or one time too few. {\small\color{kcgray}(week 8)}
\item[Reflection] Thinking back about what you did, what you learned and what you want next. {\small\color{kcgray}(week 12)}
\item[Showcase] An event where you show your work to others. {\small\color{kcgray}(week 12)}
\item[text code] Code made of typed words and symbols. {\small\color{kcgray}(week 6)}
\item[Version] A saved copy of your program at one moment, like \texttt{fish\_party\_v2.py}. {\small\color{kcgray}(week 11)}
\end{description}
\end{multicols}


\clearpage
\section{My own words}
Add new words you meet, and write the meaning in your own words.
\begin{center}\renewcommand{\arraystretch}{2.2}
\begin{tabularx}{\linewidth}{|p{4cm}|X|}
\hline
\textbf{Word} & \textbf{What it means (in my words)}\\ \hline
 & \\ \hline
 & \\ \hline
 & \\ \hline
 & \\ \hline
 & \\ \hline
 & \\ \hline
 & \\ \hline
 & \\ \hline
\end{tabularx}
\end{center}
\end{document}
